\chapter{Reconstruction and execution procedures}
\label{ch:execution}

\section{Ingest, reconcile and construct histories}

The historical procedure starts with a filing manifest. It records documents,
reporting dates, publication dates, revisions and expected statements and
notes. Tagged filings are parsed through an XBRL adapter; PDF and HTML sources
are segmented into text, tables, captions and references. The original bytes
and extraction output remain immutable. Re-running an improved extractor
creates a new extraction version with links to the same source.

For each table, the importer first resolves headers, units, signs, periods and
dimensions. It then proposes concept mappings. A deterministic validator checks
subtotal arithmetic, duplicate facts, comparative consistency and cross-note
references. Every proposed repair creates a new reading. The system may suggest
that an OCR value is 10 rather than 70, but only a reviewed correction or an
explicit alternative allows that reading into a validated model.

Reconciliation proceeds from local totals to movement schedules. A debt bridge
may contain opening debt, issuance, repayment, acquired debt, interest
accretion, currency effects, fair-value changes and closing debt. Every
component receives a movement type. An unexplained residual is stored as a
residual with a source and period; it is not automatically classified as
repayment or issuance. A model can introduce an unknown decomposition of that
residual, subject to evidence and economic constraints.

The reconstruction engine then performs the following finite procedure within
a declared model scope:
\begin{enumerate}
\item Build shared variables for observations, opening states, contract terms
and supported movement amounts. Attach original rounding domains.
\item Add identities, roll-forwards, account restrictions, contractual
relationships and cross-period links. Add each semantic reading as a condition
on the same joint model.
\item Propagate types, units, intervals and simple equalities. Reject context
mismatches before invoking a numerical solver.
\item Ask the solver for feasibility. On infeasibility, return a sufficient
conflicting subset of constraints, its source spans and model assumptions.
A solver's conflict core need not be the smallest possible explanation.
\item Retain finite alternatives and continuous solution families symbolically.
Materialise example histories only when needed for review or execution.
\item Check whether each requested simulation has sufficient economic state.
Missing maturities, reset dates or cohort composition become declared
assumptions or explicit unresolved inputs.
\end{enumerate}

Unrestricted journal enumeration would be unbounded: arbitrary offsetting
transactions can be inserted without changing reports. The model therefore
declares allowed event types, time resolution, relevant accounts, lot/cohort
resolution and limits on latent contract instances. These limits describe the
represented family; they are not evidence that all real possibilities were
searched. The compiler canonicalises syntactic duplicates and equivalences
established by its supported rewrite rules. It does not decide arbitrary
economic equivalence. Probability measures attach to declared economic
hypotheses or parameters, never to a count of encodings; adding a duplicate
representation therefore cannot create extra probability mass or coverage.

The first implementation uses symbolic linear bridges and finite structural
choices. Nonlinear economic models can be evaluated on feasible states through
plugins. A solver timeout preserves the unresolved family and returns
\code{SolverIncomplete}; it does not prune the family as impossible. Search
budgets belong to run configuration and are recorded with the results.

\section{Condition a run without rewriting history}

A run selects an evidence view and cutoff, a historical family, a forecast
starting date, scenario definitions, accounting packages, management policies,
output queries and numerical settings. A run may use a concrete history,
an exact symbolic family, a finite set of conditional histories, or weighted
draws from a declared distribution. The output records which mode was used.

Suppose no debt maturity is disclosed. An analyst may run a conditional case
that assumes all debt matures after the forecast horizon. The run is valid
under that assumption even though the historical representation remains
incomplete. The interface must display the assumption next to the result and
provide sensitivity cases when the assumption is material. It must not label
the assumed maturity as a discovered historical fact.

Management behaviour normally uses only information available at the decision
date. A perfect-foresight stress case is a separate, explicitly labelled
policy. Scenarios should share the same underlying innovations when estimating
paired differences where appropriate, with stable stream identifiers so
branch order does not change the draws. Common random numbers improve the
comparability of paired simulations; they do not imply that every branch has
the same economic event counts or shock distribution.

\section{Schedule and commit events}

An annual cash reconciliation leaves a practical question unanswered: can the
company pay when each obligation falls due? Northstar opens its projection
with cash 74.6, plans collections of 280 and supplier payments of 160. If
collections arrive first, that supplier payment can be made from cash. If
the supplier must be paid first, the company is short by 85.4. Adding the
same annual receipts and payments gives the same unconstrained year-end cash
in both orders. Only one order establishes this payment's feasibility when no
eligible financing is available.

This is a conditional timing example, not a newly inferred invoice history.
The original Northstar projection places collections before payments. The
alternative schedule shows why aggregate annual drivers need dated obligations
before they can support a liquidity conclusion. Unknown collection dates or
unknown facility access leave feasibility unresolved; neither implies default.
The proposed schedule annotation in
Section~\ref{sec:settlement-schedule-draft} makes these distinctions explicit.

An explicit borrowing decision can repair this timing shortage, but its costs
must follow through the accounts. Suppose the scenario supplies an eligible
facility and chooses a draw of 85.4 on 4 January, before paying the supplier.
Collections arrive on 5 January and the draw is repaid on 6 January. Under an
authored 5\% simple Actual/365 contract, two calendar days of interest cost
\(85.4(0.05)(2/365)=427/18250\). With the demonstration's immediate 20\% tax
method and deductible interest, annual profit, equity and closing cash each
fall by \(854/45625\). The principal inflow and repayment remain separate
financing movements even though they net to zero. The full-year comparison
in Section~\ref{sec:funded-year-authoring} accounts for those costs explicitly.

The commitment, collection dates and tax treatment are scenario assumptions.
A proposed draw is not cash already received. If it settles and a later
payment fails, the borrowing remains; refusing the payment cannot erase the
earlier event. In the conditional reference example, moving payroll before
collections produces a cash shortage of 50. Payroll remains unapplied while
the completed draw and supplier settlement remain in the state.

A report of that failed path must distinguish the available snapshot from
the requested annual statements. Cash immediately after the supplier payment
is known to be zero. Year-end cash and profit remain unresolved because the
later operations have not completed. The report retains the last successful
event, remaining obligations and known opening balances; it does not rename
the interrupted state a year-end result. The operation-driven example in
Section~\ref{sec:company-input-authoring} checks this distinction alongside
the complete cash and equity bridges.

The same distinction governs a sequence of years. An incomplete year supplies
no closing state from which the next year can start. Earlier completed years
remain available, while a revision of their assumptions requires the affected
later years to be recomputed. Section~\ref{sec:company-operation-paths}
checks these dependencies with repeated claims and asset cohorts, rather than
carrying only annual account totals.
Section~\ref{sec:company-path-authoring} then distinguishes an exchanged partial
outcome from a closing state, and a fixed collection instruction from an
explicit rule that regenerates it after a price change.

Each event carries an economic occurrence time, accrual interval when relevant,
settlement time, entity, contract or cohort identifier, currency, inputs and
provenance. Its ordering key consists of date, declared phase, dependency edges
and a deterministic tie identifier. The reference phases are: opening and
scheduled contractual changes; market observations and rate resets; business
events and decisions; accrual and measurement; settlement; period closing and
reporting. A contract can require a more specific dependency. These phases are
an engineering starting point, not a universal description of market timing.

Two events may share a tie order only if their order is either specified or
irrelevant to the result. If payment before remeasurement and payment after
remeasurement differ, an unresolved ordering is a model ambiguity or an error.
The engine does not conceal it behind alphabetical event IDs.

\begin{lstlisting}
execute_event(operation_request):
    validate exact input closure, bindings and historical admission
    validate separate posting authorization, journal and delta
    lookup stable entity/book/branch/business-event key
    if exact committed fingerprint: return original receipt
    if different committed fingerprint: return event conflict
    compare expected opening state and head generation
    atomically recheck key/head and commit journal, delta, state, receipt
    return receipt, emitted dependent events and diagnostics
\end{lstlisting}

The atomic unit includes the economic state and its accounting effects.
If posting fails, neither state changes. A repeated event ID with identical
inputs returns the existing result; the same ID with different inputs is a
conflict. Here the inputs include the request identifier, expected head and
execution settings, not only the financial payload. Section~\ref{sec:replay}
defines the stable key and full fingerprint. Consolidation runs after entity books reach the required date and
stores its adjustments separately. A failure in one scenario branch cannot
mutate another branch's state.

\section{From an economic event to permission to post}
\label{sec:posting-permission}

A delivery can change revenue, inventory and receivables at once. Before any
book changes, the method needs to know which contract was performed, which
opening state it consumes, which evidence supports that reading and which
accounting profile applies. Balance alone cannot supply those facts. AccDSL
therefore separates the economic basis, permission to use a treatment, proposed
entries and the transaction that makes them durable.

The event basis binds an immutable event to opening state, context, branch,
time view, chart, authority, method, evidence, admissions, assumptions and view.
Its business identity is the four-tuple of entity, book, branch and business
event. The branch declares whether it is synthetic or real. A synthetic teaching
profile can justify a teaching posting; it cannot authorize a real book merely
because the resulting arithmetic is identical.

The method first evaluates reporting framework, jurisdiction and dates; then
entity and transaction scope; contract and recognition facts; elections;
numerical capability; and account/context bindings. Every material finding
survives. A false required fact is ineligible, an unknown required fact is
unresolved, and an event outside the method is outside scope. A known false
finding must not disappear behind a later unknown one. Missing operative
authority never becomes source-supported permission.

The native operation \code{accounting.assess@1} remains read-only assessment.
Its \code{assess\_only} request and false posting permission cannot be promoted
by editing a flag. The separate \code{accounting.authorize@1} operation produces
a \code{PostingAuthorization} from the complete basis and a supported posting
profile. It permits posting only when the treatment is eligible, the method is
fully specified for the event and the branch mode matches the profile's
evidence scope. An unresolved or unsupported assessment is still a useful
completed assessment, but it contains no permission to post.

An authorized proposal contains both a journal and an independent economic
delta, each bound to the same basis and authorization. Journal lines name
account role, debit or credit, positive exact quantity, currency/context,
entity, book, event and method. Zero lines are omitted; a negative effect reverses
sides. The journal balances within each currency and book. The delta separately
checks quantities and contractual memory: inventory movements consume available
lots, cash follows settlement, share counts follow ownership events and debt
face amount remains distinct from carrying value.

The commit arguments have exactly four references:
\begin{lstlisting}
AccountingCommitArguments/v1 {
  basis_ref, authorization_ref, journal_ref, delta_ref
}
\end{lstlisting}
The expected opening state and branch generation appear once, in the outer
operation request. A reference to a native assessment in the authorization
field fails its reference type. False posting permission fails authorization.
An unknown posting schema returns unsupported operation with no journal commit.
All bindings are checked before looking up a prior receipt.

\begin{figure}[htbp]
\centering
\begin{tikzpicture}[node distance=5mm,>=Stealth,
 box/.style={draw,rounded corners,align=center,text width=29mm,minimum height=13mm,font=\small}]
\node[box] (basis) {Economic event\\and opening state};
\node[box,right=of basis] (auth) {Eligible method\\and permission};
\node[box,right=of auth] (proposal) {Journal and\\economic delta};
\node[box,right=of proposal] (commit) {Atomic state\\and receipt};
\draw[->] (basis)--(auth);
\draw[->] (auth)--(proposal);
\draw[->] (proposal)--(commit);
\end{tikzpicture}
\caption{Each step answers a different question. Accounting assessment may
inform the permission decision, but only the final transaction changes the book.}
\end{figure}

For a delivery, the basis identifies the exact performance obligation,
quantity, allocation and control finding. The opening state contains cumulative
revenue, billing, collection and inventory. The selected method derives the
recognition and lot movements; the report later consumes that same committed
path. Combining a journal from one method version with a delta or opening state
from another is invalid even if the totals happen to balance.

\section{A lost acknowledgement must not create a second sale}
\label{sec:replay}

Suppose a sale commits but the response is lost. The client knows the request,
but cannot tell whether the server committed it. Retrying under a new event ID
could duplicate revenue. Retrying the original request must instead return the
original immutable receipt.

The lookup key is
\[
  (\text{entity},\ \text{book},\ \text{branch},\ \text{business event}).
\]
Payload, event version, opening state and request settings belong to the complete
fingerprint, not this key. A renamed run or bumped event version cannot evade a
conflict. A separately commissioned hypothetical projection uses a new branch
namespace; reopening the same projection uses its existing one.

A correction of an already committed event is a new business event linked to
its predecessor and an admitted correction method, such as reversal and
replacement or a dated adjustment. It preserves the original receipt.
The accounting treatment must apply to the actual correction; absent treatment
authority leaves it unresolved rather than licensing a generic reversal.

Commit first validates input shape, closure, bindings and authorization in the
requested historical view. It then looks up the stable event key. An equal
fingerprint returns the old receipt with no additional commit; a different
fingerprint returns an event-identity conflict. Only an absent key proceeds to
the expected-state and generation comparison. This order is essential: after
a successful sale the current head has changed, so checking freshness before
replay would reject the very retry needed to resolve the lost response.

A new candidate is assessed and constructed without mutation. One transaction
then rechecks key absence and the expected head, inserts journal, delta, closing
state, dependencies and receipt, and advances the head. It commits all or none.
If another client has committed the identical event, return that receipt. A
different event that advanced the head makes this request stale. A changed
payload under the same key conflicts. Diagnostic or scratch records never count
as economic transitions.

A timeout consequently has three possible observations: committed, not committed
or unknown. Unknown retains the original key, fingerprint and expected head for
receipt lookup or exact retry. Refreshing the head changes the request and cannot
be presented as the same retry. Missing committed dependencies are an integrity
failure, not permission to reconstruct an approximately similar receipt.

\section{Funding, tax and other simultaneous relations}

Some financial models have circular dependencies. Interest affects taxable
profit; tax and investment affect funding needs; funding affects interest.
The reference executor first tries to resolve these through explicit timing.
For example, beginning-of-period debt can determine current interest, with
new borrowing affecting the next accrual period. That is acceptable only when
the model's timing assumption is stated and appropriate.

Where the economic model genuinely requires simultaneous decisions, the
dependency graph identifies the strongly connected component and delegates it
to a declared equation or optimisation solver. The solver must return a
solution, residuals, admissibility checks and uniqueness or multiplicity
status. A contraction iteration requires a justified convergence condition;
an arbitrary fixed iteration count is not a mathematical solution. Multiple
admissible equilibria become alternatives when material. No solution produces
a model outcome such as financing failure, or an explicit numerical failure,
depending on the cause.

Funding policy must include available facilities, limits, covenants, collateral
and allowed capital actions. The program cannot issue unlimited debt to keep
cash positive. An inability to meet payments may be a legitimate default
outcome. Default ends or changes contractual and accounting processes under
the relevant event rules; it is not quietly discarded as a bad simulation draw.

\section{Worked example: identical history, different future journals}

Return to the company in Chapter~\ref{ch:purpose}. Immediately after the
historical interest payment and closing, cash is 95, debt is 100 and equity
is $-5$. There are no other assets, liabilities or transactions. Next year's
expense is 5 in the all-fixed branch and 7 in the all-floating branch under
the rate shock. The generated entries and closing statements are:

\begin{center}
\begin{tabular}{lrr}
\toprule Output & All fixed & All floating\\\midrule
Debit interest expense / credit cash &5&7\\
Closing cash &90&88\\
Closing debt &100&100\\
Current-year profit & $-5$ & $-7$\\
Closing equity after profit transfer &$-10$&$-12$\\
Cash change &$-5$&$-7$\\\bottomrule
\end{tabular}
\end{center}
The balance-sheet identity holds in each branch: $90=100-10$ and
$88=100-12$. The demonstration reporting profile classifies interest paid as
operating cash flow. A different permitted reporting profile changes that
classification, while preserving the cash movement and source event. This
example proves the intended distinction and its arithmetic; it does not
estimate the probability that either debt structure describes a real company.

\section{Worked example: a bond has several relevant amounts}

A two-period zero-coupon bond raises cash of 100 and promises 110.25 at maturity.
Assume a 5\% effective rate per period, no transaction costs, no default, and
amortised-cost treatment under a demonstration rule. The effective-interest
schedule follows directly from $100(1.05)^2=110.25$:
\begin{center}
\begin{tabular}{lrrr}
\toprule Date & Opening liability & Interest accretion & Closing liability\\\midrule
Issue &0&0&100\\
End of period 1 &100&5&105\\
End of period 2, before payment &105&5.25&110.25\\
Redemption &110.25&0&0\\\bottomrule
\end{tabular}
\end{center}
Issuance debits cash and credits the liability by 100. Each accretion debits
interest expense and credits the liability. Redemption debits the liability
and credits cash by 110.25. Legal face amount, carrying amount and market
fair value remain separate fields. A market-rate shock may change fair value
without changing this contractual payment schedule or this demonstration
amortised-cost path.

Cash-flow reporting uses the actual payment event and its economic components;
the annual accretion is non-cash. The applicable reporting package decides how
to classify the maturity payment, rather than inferring it solely from the
liability debit. Fees, modifications, early redemption, embedded options,
hedges and foreign currency require additional event and measurement rules.

Ordinary equity issuance similarly separates cash, share quantity and legal
capital. Ten new shares at 12, with par value 1, produce cash of 120, share
capital of 10 and premium of 110 under a simple par-value profile. Shares
outstanding increase by ten units. Direct issuance costs of 3 reduce the
specified equity account and cash under the chosen illustrative treatment.
The accounting rule and the cost's eligibility must be explicit; this example
does not prescribe every jurisdiction's treatment.

\section{Reporting and explanations}

The reporting engine produces the balance sheet, income statement, comprehensive
income, cash-flow statement, changes in equity and note schedules when the
necessary modules are present. Each output carries its basis, perimeter,
period, scenario and interpretation. A comparative table cannot combine an
original prior-year figure and a restated current-year comparative without
displaying that choice.

An explanation is a dependency graph with typed edges: extracted from,
interpreted as, assumed by, constrained by, generated by, measured by,
posted by, aggregated into and supersedes. A cash-flow line must lead back
to payment events; a fair-value movement must lead back to valuation inputs;
a consolidated amount must show eliminations. A provenance path establishes
how a result was obtained. It does not by itself establish that the underlying
economic model is correct.

\section{The corporate reference package as operations}

The Northstar example defines a closed, reproducible corporate package rather
than a universal standard. Its events separate performance, invoicing and
settlement even where the example makes the first two simultaneous. Every
payload includes entity/book, currency, economic date, period, source or
assumption IDs and an information-view ID. Amounts are exact before the
declared posting quantisation. Quantities and money are checked separately.

\begin{longtable}{L{27mm}L{60mm}L{49mm}}
\toprule Event & Preconditions and state change & Posting\\\midrule\endhead
Deliver/invoice & Goods owned and available; control transferred; quantity
removed at selected cost; unconditional amount due added & Debit receivable,
credit revenue; debit cost of sales, credit inventory\\
Collect & Receipt no greater than selected unpaid receivable; reduce that
balance; excess must be separately classified & Debit cash, credit receivable\\
Receive goods & Accepted owned quantity at supported unit cost; create supplier
obligation and inventory lot & Debit inventory, credit payable\\
Pay supplier & Identified obligation and sufficient available cash; reduce
both & Debit payable, credit cash\\
Payroll & Service period and cash obligation evidenced or assumed & Debit
payroll expense, credit cash; accrual variant uses a payable\\
Acquire plant & Control obtained, cost known, in-service date/life/residual
recorded; add asset cohort & Debit plant cost, credit cash or payable\\
Depreciate & In-service asset with remaining depreciable amount; apply dated
schedule, capped by that amount & Debit depreciation, credit accumulated
depreciation\\
Interest & Contract, reset state and accrual fraction known; create interest
obligation, then settle if cash available & Debit interest, credit interest
payable; debit payable, credit cash\\
Current tax & Reference tax base positive; 20\% teaching rule; apply prepaid
tax before further payment & Debit tax expense; credit prepaid tax or payable;
settlement separately affects cash\\
Dividend & Explicit authorised distribution, sufficient funds and permitted
capital action & Debit retained earnings, credit dividend payable; then settle\\
Close & Required period operations completed; preserve original movements &
Transfer temporary income/expense balances to retained earnings\\\bottomrule
\end{longtable}

The current example combines accrual and settlement when they share a date,
but the separate obligation remains the general operation. Customer advances
are debit cash/credit contract liability and become revenue only on the
supported performance event. Inventory costing requires an admissible policy;
negative inventory and collection beyond an identified claim are rejected.
Non-cash acquisitions, write-offs, impairment and currency movements require
their own classified templates. An unknown residual never enters one of these
operations merely to force a reconciliation.

Reporting maps asset and contra-asset balances into the balance sheet and
liability/equity balances with their opposite sign. Before closing, current
profit is included once in equity; after closing it is already in retained
earnings. Income uses period movements in income/expense accounts, and OCI
uses its own designated accounts. Cash flow sums cash legs by economic
operation: Northstar's collections, suppliers, payroll, interest and taxes
are operating, plant payments investing, and debt/equity proceeds or repayments
and dividends financing. Non-cash depreciation is excluded from direct cash
flow and added back only in the indirect reconciliation. Note schedules use
the same contracts, lots and classified movements as the statements.

\section{Scenario meaning and information access}

An \code{ExogenousStress} supplies a path or joint driver law while holding the
historical evidence model fixed. \code{ObservationalConditioning} additionally
supplies actual observations and a joint probability model; the finite method
in Chapter~\ref{ch:math} evaluates the conditional distribution when its
normalising mass is positive. \code{StructuralIntervention} requires a named
structural model and replaces designated equations, retaining its exogenous
noise. An absent structural model gives \code{UnsupportedScenario}, not an
intervention inferred from an accounting identity \citep[Section 3.2.1]{pearl2009}.
The textual Core 0 \code{scenario} declaration is an exogenous stress only;
the other modes require the explicit typed API records.

At each decision date the runtime constructs an information view containing
eligible disclosures, committed state through that date, currently observed
drivers and the declared forecasting model. Future realised innovations and
future outcomes are inaccessible. A stress path may be announced in advance,
but this is an explicit visibility policy. Otherwise management sees only its
current and past values. A bank's allowance at the opening date may use a
forecast distribution of defaults; it must not inspect the outer simulation's
future realised default. Inner valuation randomness has a separate stream
domain from outer realised outcomes.

Funding follows a declared waterfall: available cash, committed eligible
facilities up to their limits, and separately authorised capital actions.
Each action checks collateral, covenants and settlement timing. If an amount
remains due, record an economic funding failure and the unpaid obligation.
If required valuation or accounting is unsupported, stop the affected
dependencies as incomplete. The branch manifest preserves the last committed
state, completed outputs, unavailable outputs and reason. These distinctions
let an analyst see a meaningful bad outcome without confusing it with broken
computation.
